%% ══════════════════════════════════════════════════════════════
%%  CHAPTER 1 — SYSTEM OVERVIEW
%% ══════════════════════════════════════════════════════════════
\chapter{System Overview}
\label{ch:overview}

\section{Introduction}
\label{sec:overview:intro}

\autoscholar is an autonomous scientific research engine that generates publication-quality academic papers in quantitative finance, econometrics, statistics, machine learning, and artificial intelligence---entirely without human intervention.
The system represents a fundamental departure from traditional AI-assisted writing tools: rather than augmenting a human researcher, \autoscholar \emph{replaces} the entire research pipeline, from literature review through data collection, statistical analysis, and manuscript preparation.

At its core, the platform leverages \claude---Anthropic's flagship large language model with a one-million-token context window---as the cognitive backbone of a multi-agent architecture.
Five specialized sub-agents, each named after a foundational figure in mathematics and science, collaborate through a carefully orchestrated pipeline to produce papers that target the quality standards of top-tier academic journals such as the \emph{Journal of Finance}, \emph{Journal of Financial Economics}, and the \emph{Review of Financial Studies}.

The system operates continuously.
On a configurable cadence (default: one paper per hour), \autoscholar autonomously selects a research topic from a bank of over 1,000 pre-constructed prompts spanning 24 research domains and 24 methodological approaches, executes the full research pipeline, validates the output through AI-powered peer review, and publishes the result to its self-hosted autonomous journal.
Each paper receives a persistent identifier (the AutoScholar Identifier, or \gls{asid}), is stored with full \LaTeX{} source, BibTeX references, Python analysis scripts, and raw data---packaged as a reproducible research archive.

\section{Motivation and Vision}
\label{sec:overview:motivation}

The academic research process is notoriously slow, expensive, and bottlenecked by human bandwidth.
A typical empirical finance paper requires 6--18 months from conception to submission, involving literature review (weeks), data acquisition and cleaning (weeks to months), econometric modeling (months), writing (months), and revision cycles.
\autoscholar compresses this timeline to approximately 45--90 minutes per paper by automating every stage.

The vision is not to produce perfect papers---many generated papers will contain limitations that a human reviewer would flag---but to demonstrate that the \emph{entire pipeline} of empirical research can be automated end-to-end, creating a foundation for:

\begin{enumerate}[label=(\roman*)]
  \item \textbf{Rapid hypothesis exploration.} Researchers can use \autoscholar to quickly test whether a research direction is viable before investing months of manual effort.
  \item \textbf{Reproducible research at scale.} Every paper is packaged with its complete computational pipeline, enabling verification and extension.
  \item \textbf{Continuous scientific production.} The autonomous journal model demonstrates that research output can be ``always on,'' analogous to continuous integration in software engineering.
  \item \textbf{Democratization of quantitative research.} By encoding expert-level econometric and statistical knowledge into the system prompt and tool definitions, \autoscholar makes sophisticated quantitative methods accessible to researchers without deep programming expertise.
\end{enumerate}

\section{Technology Stack}
\label{sec:overview:stack}

\autoscholar is a full-stack TypeScript/Python application with the following major technology choices:

\begin{table}[H]
\centering
\caption{Core Technology Stack}
\label{tab:tech-stack}
\begin{tabularx}{\textwidth}{lX}
\toprule
\textbf{Layer} & \textbf{Technologies} \\
\midrule
AI Engine & Claude Opus 4.6 (1M context), Claude Haiku 4.5 (vision), OpenRouter (400+ models) \\
Backend & Node.js, Express 4.21, TypeScript 5.6 \\
Frontend & React 18.3, Vite 5, TailwindCSS 3, TanStack React Query, Framer Motion \\
Database & SQLite (better-sqlite3) / PostgreSQL (Neon Serverless), Drizzle ORM \\
Python Runtime & Python 3.x sandbox with NumPy, Pandas, SciPy, Statsmodels, Scikit-learn, XGBoost, LightGBM, ARCH, CVXPY, Matplotlib \\
Document Engine & pdflatex, bibtex, 50+ \LaTeX{} packages; Puppeteer (HTML$\to$PDF); Pandoc (DOCX) \\
Financial Data & FMP (230+ endpoints), EODHD, Financial Data API (real options) \\
Macroeconomic Data & World Bank (16,000+ indicators), IMF (DataMapper + SDMX 3.0), OECD (1,475+ dataflows), FRED (840,000+ series) \\
Academic Data & Semantic Scholar, OpenAlex, arXiv, CrossRef, Google Scholar (SerpAPI) \\
Web Intelligence & Tavily (5 tools), Firecrawl (7 tools), Exa (2 tools), SerpAPI (11 verticals) \\
Prediction Markets & Polymarket (Gamma, Data, CLOB APIs) \\
\bottomrule
\end{tabularx}
\end{table}

\section{Scale and Metrics}
\label{sec:overview:metrics}

The codebase comprises approximately 40,000+ lines of source code across 155+ files, exposing 270+ HTTP API endpoints and 213 Claude-callable tool definitions.
Key quantitative characteristics include:

\begin{itemize}
  \item \textbf{Tool definitions:} 213 tools across 7 categories (financial, academic, web, macro, Python, research paper, prediction markets), totaling $\sim$6,100 lines of JSON schema definitions.
  \item \textbf{System prompt:} 1,798 lines defining the complete research pipeline, sub-agent protocols, data fallback chains, and academic writing standards.
  \item \textbf{Prompt bank:} 1,000+ deterministically generated research prompts spanning 24 domains $\times$ 24 methods $\times$ 14 topics per domain.
  \item \textbf{Python execution environment:} 17 pre-installed scientific computing libraries, 5 pre-loaded API clients, 4 embedded databases (SDC Platinum IPOs, US Housing via DuckDB), and 8 file export functions.
  \item \textbf{Financial data coverage:} 230+ FMP endpoints, EODHD historical data to IPO date, real-time options with Greeks from Financial Data API, and 840,000+ FRED macroeconomic series.
  \item \textbf{Tool execution depth:} Up to 180 recursive tool calls per auto-paper run (45 for interactive chat).
\end{itemize}

\section{Repository Structure}
\label{sec:overview:repo}

The repository follows a monorepo structure with clear separation between client, server, and shared code:

\begin{figure}[H]
\centering
\begin{forest}
  for tree={
    font=\ttfamily\small,
    grow'=0,
    child anchor=west,
    parent anchor=south,
    anchor=west,
    calign=first,
    edge path={
      \noexpand\path [draw, \forestoption{edge}]
      (!u.south west) +(7.5pt,0) |- node[fill,inner sep=1.25pt] {} (.child anchor)\forestoption{edge label};
    },
    before typesetting nodes={
      if n=1
        {insert before={[,phantom]}}
        {}
    },
    fit=band,
    before computing xy={l=15pt},
  }
[autoscholar/
  [client/
    [src/
      [components/ \textnormal{\textit{(chat, UI, streaming)}}]
      [hooks/ \textnormal{\textit{(useChatSession, useChatMutation, useAnalytics)}}]
      [pages/ \textnormal{\textit{(home, admin, showcase, shared-report, ...)}}]
      [utils/ \textnormal{\textit{(pdfGenerator)}}]
      [lib/ \textnormal{\textit{(queryClient)}}]
    ]
  ]
  [server/
    [routes/ \textnormal{\textit{(chat, admin, reports/, api/v1/)}}]
    [services/
      [claude/ \textnormal{\textit{(systemPrompt, streaming, tools/)}}]
      [python/ \textnormal{\textit{(executor, GARCH, VaR, Monte Carlo, ...)}}]
      [latex/ \textnormal{\textit{(compiler, generator, sanitizer, patch)}}]
      [governor/ \textnormal{\textit{(engine, profiles, decision log)}}]
      [fmp/ \textnormal{\textit{(client, company, quotes, ...)}}]
      [pdf/ \textnormal{\textit{(Puppeteer, Pandoc, HTML export)}}]
      [prompt-bank/ \textnormal{\textit{(domains, topics, index)}}]
    ]
    [storage/ \textnormal{\textit{(paper, message, session, analytics, admin, user)}}]
  ]
  [shared/
    [schema.ts \textnormal{\textit{(PostgreSQL Drizzle schema)}}]
    [schema-sqlite.ts \textnormal{\textit{(SQLite Drizzle schema)}}]
    [types.ts \textnormal{\textit{(shared API types)}}]
  ]
  [downloads/ \textnormal{\textit{(generated CSV, XLSX, JSON, ZIP files)}}]
  [logs/ \textnormal{\textit{(daily log files)}}]
  [public/ \textnormal{\textit{(static assets, generated plots)}}]
]
\end{forest}
\caption{Top-level repository structure of the \autoscholar codebase.}
\label{fig:repo-structure}
\end{figure}

\section{Key Design Decisions}
\label{sec:overview:design}

Several architectural decisions shape the character of the system:

\begin{description}[style=nextline]
  \item[Monolithic streaming loop.] Rather than using a message-queue or actor framework, the entire research pipeline executes within a single HTTP streaming connection.
  Claude issues tool calls, the server executes them synchronously, returns results, and Claude continues---all within one SSE stream that can run for 45--90 minutes.
  This simplifies state management but creates challenges for resilience (addressed by the checkpoint system).

  \item[Claude as the cognitive core.] All reasoning, planning, and writing is performed by \claude.
  The sub-agents (\gauss, \turing, \newton, \fisher) are not separate model instances but \emph{conceptual roles} described in the system prompt.
  Claude switches between these roles based on the stage of the pipeline.
  The only exception is the \governor, which is a pure rule-based engine with zero LLM calls.

  \item[TypeScript--Python bridge pattern.] Computationally intensive operations (GARCH estimation, portfolio optimization, Monte Carlo simulation, custom analysis) run in Python subprocesses spawned by TypeScript wrappers.
  Each TypeScript module validates inputs, spawns the Python script, parses structured JSON output, and handles errors.
  This pattern appears in 10+ service pairs across the codebase.

  \item[Dual database support.] The system supports both SQLite (for development and self-hosted deployments) and PostgreSQL (for production via Neon Serverless).
  Drizzle ORM provides a unified query interface, with separate schema definitions for each backend.

  \item[Tool-rich agent paradigm.] With 213 tool definitions, \autoscholar takes the ``tool-rich'' approach to AI agents rather than the ``few-tools'' approach.
  This gives Claude maximum flexibility but requires careful prompt engineering to guide tool selection and prevent misuse.
\end{description}

\section{Document Organization}
\label{sec:overview:organization}

The remainder of this document is organized as follows:

\begin{description}
  \item[\Cref{ch:architecture}] presents the complete system architecture, including component diagrams, dependency graphs, and the execution model.
  \item[\Cref{ch:agents}] provides deep dives into each of the five sub-agents plus the \governor and \popper validator.
  \item[\Cref{ch:pipeline}] reconstructs the full paper generation pipeline from trigger to publication.
  \item[\Cref{ch:journal}] describes the autonomous journal system, including scheduling, quality control, and publication mechanics.
  \item[\Cref{ch:data-layer}] details the data acquisition, storage, and management infrastructure.
  \item[\Cref{ch:modeling}] examines the modeling and scientific computation layer.
  \item[\Cref{ch:paper-engine}] explains the \LaTeX{} paper generation and compilation engine.
  \item[\Cref{ch:memory}] covers the memory, state, and deduplication systems.
  \item[\Cref{ch:failure}] analyzes failure handling, retry logic, and recovery mechanisms.
  \item[\Cref{ch:performance}] discusses performance, scalability, and resource constraints.
  \item[\Cref{ch:design-patterns}] identifies architectural patterns and design philosophy.
  \item[\Cref{ch:weaknesses}] provides a critical evaluation of system weaknesses.
  \item[\Cref{ch:improvements}] proposes improvements for the next generation.
  \item[\Cref{ch:visualization}] contains system diagrams and pipeline visualizations.
  \item[\Cref{ch:frontend}] documents the frontend architecture.
  \item[\Cref{ch:api}] specifies the public REST API.
\end{description}
